\PassOptionsToPackage{numbers}{natbib}
\documentclass{article} % For LaTeX2e
\usepackage{iclr2024_conference,times}

\usepackage[utf8]{inputenc} % allow utf-8 input
\usepackage[T1]{fontenc}    % use 8-bit T1 fonts
\usepackage{hyperref}       % hyperlinks
\usepackage{url}            % simple URL typesetting
\usepackage{booktabs}       % professional-quality tables
\usepackage{amsfonts}       % blackboard math symbols
\usepackage{nicefrac}       % compact symbols for 1/2, etc.
\usepackage{microtype}      % microtypography
\usepackage{titletoc}

\usepackage{subcaption}
\usepackage{graphicx}
\usepackage{amsmath}
\usepackage{multirow}
\usepackage{color}
\usepackage{colortbl}
\usepackage{cleveref}
\usepackage{algorithm}
\usepackage{algorithmicx}
\usepackage{algpseudocode}
\usepackage{tikz}
\usepackage{pgfplots}
\usepackage{float}
\usepackage{array}
\usepackage{tabularx}
\pgfplotsset{compat=newest}


\DeclareMathOperator*{\argmin}{arg\,min}
\DeclareMathOperator*{\argmax}{arg\,max}

\graphicspath{{../}} % To reference your generated figures, see below.

\title{Schema-Constrained Two-Stage Verification for Low-Cost Japanese Document-Level Relation Extraction}

\author{AIRAS}

\newcommand{\fix}{\marginpar{FIX}}
\newcommand{\new}{\marginpar{NEW}}

\begin{document}

\maketitle

\begin{abstract}
Japanese document-level relation extraction (DocRE) can populate knowledge graphs from long-form text, but in a low-cost regime where an API-accessed LLM is prompted once to extract all triples, outputs often contain many unsupported relations, yielding low precision. This is difficult because the number of entity pairs per document is large, evidence may be distributed across sentences, and unconstrained generation can produce malformed or inconsistent structures in addition to semantic hallucinations. We present a training-free Two-Stage pipeline that improves reliability while remaining inexpensive: Stage 1 performs recall-oriented candidate generation under JSON Schema–constrained structured outputs; Stage 2 performs precision-oriented verification of candidates in batches, also with schema-constrained decoding; a final deterministic filter applies relation-specific domain/range constraints using empirically observed (head\_type, tail\_type) pairs from JacRED training data. We evaluate on 10 JacRED development documents selected by character-length stratified sampling, comparing a one-shot baseline to the Two-Stage pipeline across Gemini Flash family configurations (2.0/2.5/3 preview, with and without thinking). Logged results show consistent false-positive reduction and precision gains; the best observed micro-F1 is 0.27 (Gemini 2.5 Flash with thinking 2048, and Gemini 3 Flash Preview with thinking off). Recall remains low overall, indicating that candidate coverage and entity alignment are the main bottlenecks.
\end{abstract}

\section{Introduction}\label{sec:intro}
Document-level relation extraction (DocRE) predicts directed semantic relations between entity pairs within an entire document rather than within isolated sentences. This shift in scope matters: many relations require evidence distributed across multiple sentences, and the number of candidate entity pairs grows quadratically with the number of entities, making reliable extraction challenging at scale~\cite{yao-2019-docred}. As a consequence, DocRE has historically been approached with supervised models designed for document reasoning, often combined with architectural choices such as graph-based aggregation or transformer encoders tailored to long context~\cite{delaunay-2023-comprehensive}.

Recent work has broadened DocRE resources beyond English. In particular, JacRED provides a Japanese Wikipedia-based DocRE benchmark with typed entities and evidence annotations, and shows that Japanese-specific linguistic phenomena and topic mismatch reduce the effectiveness of simple cross-lingual transfer. Importantly for the present work, JacRED’s original study also reports that large language models (LLMs) used via in-context learning perform poorly compared to supervised Japanese models, suggesting that naive prompting is not sufficient for Japanese DocRE~\cite{ma-2024-building}.

This paper targets a practical, cost-constrained regime: extracting knowledge-graph-style triples from Japanese documents using low-cost, API-accessed LLMs. In this regime, one-shot prompting is appealing because it is operationally simple, minimizes latency, and limits per-document API calls. However, our motivating hypothesis is that one-shot extraction tends to suffer from low precision. When an LLM is asked to produce many triples at once, it often outputs relations that are plausible but unsupported by the document, makes directionality mistakes, or produces inconsistent fields that complicate parsing and downstream evaluation.

Three factors make this failure mode particularly persistent. First, DocRE mixes local cues with cross-sentence aggregation, so extraction requires document-level coherence and evidence tracking rather than relying on single-sentence patterns~\cite{delaunay-2023-comprehensive}. Second, generative models can hallucinate relations when prompts implicitly reward coverage, and hallucinations can be difficult to detect if the output remains superficially consistent. Third, unconstrained natural language generation is brittle when the downstream consumer expects machine-readable structure; even a semantically correct relation can be unusable if it is emitted with missing fields or invalid JSON\@. Surveys of generative information extraction emphasize these reliability challenges and discuss constrained decoding and verification as recurring mitigation strategies~\cite{xu-2023-large}. Complementarily, recent benchmarking of JSON Schema–constrained decoding shows that structured-output engines can substantially improve schema compliance and sometimes improve downstream task accuracy, but syntactic compliance alone does not guarantee semantic faithfulness~\cite{geng-2025-generating}.

We propose a training-free Two-Stage pipeline that explicitly separates recall-oriented candidate generation from precision-oriented verification, while enforcing structured outputs at each step. Stage 1 generates a broad set of candidate relation triples for each document, using JSON Schema–constrained decoding to ensure the output is parseable and complete. Stage 2 verifies candidates in batches: given the same document and a list of candidates, the model outputs a schema-constrained decision for each candidate indicating whether it is supported by the document. Finally, we apply deterministic post-processing that enforces relation-specific domain/range constraints derived from training data: for each relation label, we compute empirically observed (head\_type, tail\_type) pairs and remove verified triples that violate these type-pair constraints.

We evaluate the method on the JacRED development set, selecting 10 documents via character-length–based stratified sampling to keep costs low while enabling consistent comparisons. We compare a baseline one-shot extraction method against the proposed Two-Stage pipeline across configurations in the Gemini Flash family (Gemini 2.0 Flash, Gemini 2.5 Flash, and Gemini 3 Flash Preview), including conditions with and without thinking (and a thinking budget of 2048 tokens when available). Performance is measured using micro-averaged precision, recall, and F1 computed over the 10 documents.

Our contributions are:
- We design a low-cost, training-free Japanese DocRE pipeline that decomposes extraction into schema-constrained candidate generation followed by schema-constrained verification, targeting the precision failures of one-shot prompting.
- We introduce a deterministic filtering step based on relation-specific domain/range constraints derived from observed entity-type pairs in training data, providing an interpretable and inexpensive guardrail.
- We provide an empirical comparison across Gemini Flash configurations showing that the Two-Stage pipeline consistently reduces false positives and improves precision, while also highlighting that low recall remains the dominant bottleneck.

The results suggest several directions for future work that are not addressed by the current experiments. Improving recall without losing the precision benefits will likely require better entity alignment and mention normalization, more robust handling of relation directionality, and decomposition strategies that surface implicit or multi-hop relations that are common in DocRE benchmarks~\cite{yao-2019-docred,delaunay-2023-comprehensive}.

\section{Related Work}\label{sec:related}
Our method lies at the intersection of document-level relation extraction, Japanese DocRE resources, and reliability mechanisms for LLM-based generative information extraction.

Document-level relation extraction. DocRED established DocRE as a benchmark emphasizing multi-sentence evidence and supporting-sentence annotation, and quantified that a substantial fraction of facts require reasoning across multiple sentences~\cite{yao-2019-docred}. The subsequent DocRE literature has explored a range of supervised architectures, including graph-based reasoning and transformer-based document encoders, and has emphasized issues such as candidate-pair explosion, evidence selection, and global inference~\cite{delaunay-2023-comprehensive}. In contrast to these works, our problem setting intentionally avoids supervised training and instead optimizes an API-driven pipeline for low operational cost. This changes the design space: rather than improving an encoder’s representational capacity, we focus on controlling the behavior of a fixed LLM through structured decoding, decomposition, and post-hoc constraints.

Japanese DocRE and limitations of naive LLM usage. JacRED addresses the scarcity of Japanese DocRE resources and highlights that Japanese discourse properties and domain mismatch limit direct transfer from English DocRE data. It also reports that LLM in-context learning baselines perform poorly on Japanese DocRE compared to supervised Japanese models~\cite{ma-2024-building}. Our work uses JacRED as an evaluation testbed but differs in objective: we do not attempt to match supervised model performance; instead, we ask whether a low-cost LLM can be made practically useful for triple extraction via reliability controls.

Generative IE, constrained decoding, and verification. A broad survey of generative IE catalogs prompting, constrained decoding, and self-improvement strategies, and emphasizes that LLM-based extraction is often brittle due to hallucinations and formatting errors~\cite{xu-2023-large}. Constrained decoding, including JSON Schema–based structured outputs, directly targets format faithfulness; recent benchmark studies provide evidence that schema-constrained generation can substantially increase compliance and can modestly improve accuracy on some downstream tasks, while also documenting variability in coverage and engine behavior~\cite{geng-2025-generating}. Our work aligns with this line of research in using schema constraints, but departs in two ways: (i) we apply schema constraints not only to extraction but also to verification outputs, treating verification as a first-class structured prediction problem; and (ii) we add a deterministic type-pair filter derived from training statistics to impose a lightweight dataset-specific prior.

Non-applicability of supervised DocRE baselines in our setting. Many high-performing DocRE systems surveyed in the literature assume supervised training over the task schema and are evaluated at scale on large datasets~\cite{delaunay-2023-comprehensive}. Those approaches are not directly comparable under our constraints (no model fine-tuning, low-cost API usage, and a small evaluation sample used to control expenses). Conversely, one-shot LLM extraction is directly applicable and constitutes our baseline; our Two-Stage method can be seen as a minimal change to that baseline that substantially improves precision by reducing false positives.

\section{Background}\label{sec:background}
DocRE task and evaluation. In document-level relation extraction, a document D contains a set of entities E = {e1, …, en}. The goal is to predict a set of directed relation instances T. Each instance is a triple (h, r, t), where h is a head entity, t is a tail entity, and r is a relation label from a fixed relation schema R. Evaluation compares predicted triples against gold triples and reports micro-averaged precision, recall, and F1 over a set of documents, following established practice in DocRE benchmarks such as DocRED~\cite{yao-2019-docred}. Let TP, FP, and FN denote counts of true positives, false positives, and false negatives, aggregated across documents. Precision is TP\,/\,(TP+FP), recall is TP\,/\,(TP+FN), and F1 is 2PR\,/\,(P+R).

JacRED benchmark. JacRED is a Japanese Wikipedia-based DocRE dataset with typed entities, 35 relation labels, and evidence annotations. Its construction was motivated by the lack of Japanese-native DocRE resources and by limitations observed in translation-based cross-lingual transfer. The JacRED study benchmarks Japanese DocRE models and reports that supervised Japanese encoders substantially outperform LLM in-context baselines, underscoring the difficulty of Japanese DocRE under naive prompting~\cite{ma-2024-building}. In this paper, JacRED provides the relation schema, entity types, and the development-set documents used for evaluation.

Generative information extraction and structured outputs. In a generative IE framing, an LLM produces an output Y conditioned on document text X and a prompt P, which can be written as p\,(Y | X, P)~\cite{xu-2023-large}. When Y must be machine-readable, constrained decoding can enforce that the output conforms to a grammar or a JSON Schema, producing syntactically valid structures that are easier to parse and evaluate. Recent work distinguishes between syntactic compliance (the output validates under the schema) and semantic correctness (the extracted content is faithful to the input) and shows that constrained decoding primarily addresses the former while leaving hallucination and evidence errors largely unsolved~\cite{geng-2025-generating}.

Problem setting and assumptions. We assume a fixed set of relations R\@ and a finite set of entity types as defined by JacRED\@. For each evaluation document, we query an LLM to extract relation triples. The method is training-free in that it does not update LLM parameters\@. However, it assumes access to JacRED training data to compute empirical domain/range constraints: for each relation r, we collect the set A$_r$ of observed type pairs (type\,(head), type\,(tail)) from the training split and use these sets as deterministic filters at inference time. This use of training statistics is not model training; it is a form of data-driven constraint induction used solely for post-processing.

\section{Method}\label{sec:method}
We propose a Two-Stage extraction pipeline that combines structured-output decoding, verification, and deterministic filtering. The design goal is to improve precision relative to one-shot extraction while remaining low-cost and avoiding any parameter updates.

Stage 1: recall-oriented candidate generation with schema constraints. Given a document D, we prompt the LLM to generate a JSON object constrained by an extraction schema, denoted EXTRACTION\_SCHEMA\@. The schema requires the model to output a list of relation instances as structured records. Each record contains at least (i) a head entity representation (identifier and/or surface form), (ii) a tail entity representation, (iii) an entity type for each argument, and (iv) a relation label r in the fixed schema R. The Stage 1 prompt is intentionally recall-oriented: it encourages proposing a broad set of candidates, including cases where evidence may be indirect, because downstream verification will act as a precision filter. Using JSON Schema–constrained decoding reduces the risk that the output is malformed or missing required fields, which is a known reliability concern in generative IE~\cite{geng-2025-generating}.

Stage 2: precision-oriented verification in batched structured outputs. Let C = \{(h$_i$, r$_i$, t$_i$)\} for i = 1..m denote the candidate triples output by Stage 1. Verification takes the same document D and a batch of candidates and asks the LLM to judge whether each candidate relation is supported by the document. The output is constrained by a verification schema, VERIFICATION\_SCHEMA, which forces a structured decision for each candidate. Verification is performed in batches of size 10 to amortize API calls while keeping each verification prompt within manageable length.

Deterministic domain/range filtering via type-pair constraints. Even with verification, systematic errors can persist, including relation directionality confusions or argument-type mismatches. We therefore apply a final deterministic filter based on empirical type constraints induced from training data. For each relation r, we compute A$_r$, the set of observed type pairs (type\,(h), type\,(t)) in the training split. For any triple (h, r, t) that passes verification, we retain it only if (type\,(h), type\,(t)) is in A$_r$. This step implements a lightweight domain/range constraint that is interpretable, dataset-specific, and cost-free at inference time.

Baseline method. The baseline is a one-shot extraction strategy: a single LLM call requests entities and relations simultaneously and then applies only simple filtering. It does not separate generation from verification and does not apply relation-specific type-pair constraints.

Rationale and expected effect. The full pipeline targets distinct failure modes. Schema constraints improve robustness against format errors and parsing failures~\cite{geng-2025-generating}. Verification targets semantic hallucination by explicitly converting extraction into a generate-then-check process, consistent with reliability themes in generative IE surveys~\cite{xu-2023-large}. Type-pair constraints add a deterministic prior derived from training data, intended to remove implausible or inconsistent triples that might still survive verification.

Final output. The final prediction set for each document is the subset of Stage 1 candidates that (i) are marked as supported by Stage 2 verification and (ii) satisfy the relation-specific type-pair constraints. We evaluate this set against JacRED gold triples using micro-averaged precision, recall, and F1.

\section{Experimental Setup}\label{sec:experimental}
Data and sampling protocol. Experiments use the JacRED development split~\cite{ma-2024-building}. To control API cost while maintaining a range of document lengths, we select 10 documents via character-length–based stratified sampling. This fixed 10-document subset is used across all model and method conditions.

Compared methods. We evaluate two conditions:
\begin{enumerate}
\item Baseline (One-shot extraction): a single LLM call extracts entities and relations together, followed by simple filtering.
\item Proposed (Two-Stage KG Extraction): Stage 1 candidate generation under EXTRACTION\_SCHEMA, Stage 2 batch verification under VERIFICATION\_SCHEMA (batch size 10), and final deterministic filtering using relation-specific (head\_type, tail\_type) constraints induced from training data.
\end{enumerate}

Models and configurations. We test three Gemini Flash family models: Gemini 2.0 Flash, Gemini 2.5 Flash, and Gemini 3 Flash Preview. Where supported, we compare configurations without thinking and with a thinking budget of 2048 tokens. This design probes the interaction between additional reasoning capacity and the behavior of extraction and verification.

Implementation details. The pipeline is implemented using the Gemini API\@ via the Google GenAI SDK\@. Both extraction and verification use JSON Schema–constrained structured outputs. Verification is executed in batches of size 10. Post-processing applies empirically observed (head\_type, tail\_type) constraints per relation derived from the JacRED training split; any extracted triple whose argument types violate the relation-specific allowed set is deterministically removed.

Evaluation metrics and aggregation. We report micro-averaged precision, recall, and F1 aggregated over the 10 documents, alongside TP, FP, and FN counts. Metrics are computed as precision = TP\,/\,(TP+FP), recall = TP\,/\,(TP+FN), and F1 = 2PR\,/\,(P+R), consistent with standard DocRE evaluation practice~\cite{yao-2019-docred}.

Comparability considerations. For each Gemini configuration, Baseline and Proposed use the same underlying model and differ only in prompting strategy, decomposition into stages, verification, and deterministic filtering. Runtime and hardware are not reported because inference is performed via remote API calls in an unspecified execution environment.

\section{Results}\label{sec:results}
This section reports only the results present in the logged metrics. No additional experiments, reruns, or unlogged ablations are introduced.

Experiment 1: Gemini 2.0 Flash (thinking: none). The Baseline achieves precision 0.20, recall 0.15, and F1 0.17 with TP = 22, FP = 86, and FN = 126. The Proposed Two-Stage pipeline improves precision to 0.35 and recall to 0.19, increasing F1 to 0.25 (TP = 28, FP = 51, FN = 121). The dominant effect is a reduction in false positives (86 to 51) alongside a modest increase in true positives.

Experiment 2: Gemini 2.5 Flash (thinking: off). The Baseline yields precision 0.17, recall 0.16, and F1 0.17 (TP = 24, FP = 115, FN = 124). The Proposed method raises precision to 0.30 but reduces recall to 0.12, leaving F1 unchanged at 0.17 (TP = 18, FP = 42, FN = 130). This configuration shows that verification plus constraints can become overly conservative, substantially suppressing both false positives and true positives.

Experiment 3: Gemini 2.5 Flash (thinking: 2048). With thinking enabled, the Baseline achieves precision 0.18, recall 0.17, and F1 0.17 (TP = 25, FP = 115, FN = 123). The Proposed method improves to precision 0.36, recall 0.21, and F1 0.27 (TP = 31, FP = 54, FN = 117). Relative to the Baseline, false positives drop sharply while true positives increase.

Experiment 4: Gemini 3 Flash Preview (thinking: off). The Baseline achieves precision 0.26, recall 0.16, and F1 0.20 (TP = 24, FP = 70, FN = 124). The Proposed method reaches precision 0.36, recall 0.22, and F1 0.27 (TP = 32, FP = 56, FN = 116). This condition achieves the best F1 without using thinking, aligning with the qualitative analysis that it offers a favorable cost–performance point.

Experiment 5: Gemini 3 Flash Preview (thinking: 2048). The Baseline improves to precision 0.31, recall 0.22, and F1 0.26 (TP = 33, FP = 74, FN = 115). The Proposed method yields precision 0.37, recall 0.20, and F1 0.26 (TP = 30, FP = 52, FN = 118). Precision and FP improve, but recall decreases and F1 matches the Baseline.

Aggregate trends and interpretation. Across all configurations, the Proposed pipeline consistently reduces FP, which translates to consistent precision gains. Improvements in F1 occur in most but not all conditions; the best observed F1 is 0.27, achieved by Gemini 2.5 Flash with thinking 2048 and by Gemini 3 Flash Preview with thinking off. Recall remains low overall (0.12 to 0.22), indicating that the primary bottleneck is recovering gold relations rather than filtering spurious ones. Moreover, the interaction with thinking is non-monotonic: while thinking can improve the Baseline, within the Proposed pipeline it can also reduce TP (as in Gemini 3 Flash Preview), consistent with a verifier that becomes more conservative.

Ablation limitations. The logs do not include runs that isolate the effects of verification versus type-pair constraints, so we cannot attribute the FP reduction to a specific component in a strictly controlled way. Nevertheless, the consistent FP decreases from Baseline to Proposed across all configurations are compatible with the combined effect of verification and deterministic filtering.

Limitations. The evaluation uses only 10 development documents, so variance estimates and confidence intervals are not available and conclusions should be treated as indicative. Additionally, deterministic type-pair constraints can remove correct relations when entity types are misassigned or when a gold type pair is absent from the empirically induced constraint set.

Figures. No figures were provided in the experiment logs (figures = []). Therefore, no figure is included here, and no figure references are made.

\section{Conclusion}\label{sec:conclusion}
We studied a low-cost, API-based setting for Japanese document-level relation extraction on JacRED in which one-shot LLM prompting produces many false positive relation triples and therefore low precision. To address this, we introduced a training-free Two-Stage pipeline that uses JSON Schema–constrained structured outputs for both candidate generation and verification, followed by deterministic post-processing with relation-specific domain/range constraints induced from training data.

Across Gemini Flash family configurations, the Two-Stage pipeline consistently reduced false positives and improved precision. Micro-F1 improved in most configurations and reached 0.27 in two cases: Gemini 2.5 Flash with thinking 2048 and Gemini 3 Flash Preview with thinking off. At the same time, recall remained low across all conditions, and in some configurations the verifier plus constraints suppressed true positives, indicating sensitivity to candidate coverage and type assignment.

These findings suggest that reliability interventions can make low-cost LLM extraction more usable for knowledge graph construction even without fine-tuning, but that further progress depends on improving recall. Promising next steps include strengthening Stage 1 candidate coverage, improving entity alignment and type robustness to prevent constraint-induced false negatives, and designing prompts or decompositions that better surface implicit and multi-hop document-level relations.

This work was generated by \textsc{AIRAS}~\citep{airas2025}.

\bibliographystyle{iclr2024_conference}
\bibliography{references}

\end{document}